\subsection{Two meanings joined by one theorem}

\begin{corebox}[First 45 minutes: antiderivatives and small contributions]
Begin by reversing elementary derivatives, then introduce Riemann sums as a model of accumulation. The two ideas should remain visibly different: one asks for a function, the other for a total over an interval.
\end{corebox}

\begin{corebox}[Second 45 minutes: the fundamental bridge and applications]
The fundamental theorem explains why a total defined by limiting sums can be computed from an antiderivative. Area, average value, displacement, work, and volume are then read as versions of one principle: integrate a local contribution over the region where it acts.
\end{corebox}

Substitution and integration by parts remain elementary reverse rules. They support the ideas of the chapter; they are not the chapter's main subject.
